% ch3_a 第3章 §3.1–§3.3.2开头 (书页 68–78 / p083–p093)
% 第3章 INTERACTING BOSON SYSTEMS
\chapter{相互作用玻色子系统}

玻色系统是凝聚态物理中最简单的系统。它描述了五花八门的物理现象，例如超流、磁性、晶体的形成，等等。它也是研究相与相变的一个模型系统。

\section{自由玻色系统与二次量子化}

\begin{keypoints}
\item 二次量子化是量子系统的一种算符描述。在二次量子化中，多体系统可以表述为一个场论。
\end{keypoints}

第一次量子化（first quantization）是利用波函数对量子系统的描述；二次量子化（second quantization）则是利用算符对量子系统的另一种描述。在二次量子化中，我们无需显式地写出波函数。例如，在谐振子的算符描述中，基态是通过算符定义的：$\hat{a}|0\rangle=0$。

现在让我们考虑一个自由的 $d$ 维玻色系统（置于体积 $\mathcal{V}=L^d$ 的盒子中），其中含有 $N$ 个全同玻色子。在第一次量子化中，系统的波函数是 $N$ 个变量的对称函数。这些 $N$ 玻色子态构成一个希尔伯特空间，记作 $\mathcal{H}_N$。为了得到玻色系统的二次量子化描述，并避免书写复杂的 $N$ 变量对称函数，我们把具有各种不同玻色子数的希尔伯特空间全部合并起来，构成一个总希尔伯特空间
\begin{equation*}
\mathcal{H}=\mathcal{H}_0\oplus\mathcal{H}_1\oplus\mathcal{H}_2\oplus\cdots\oplus\mathcal{H}_N\oplus\cdots
\end{equation*}
总希尔伯特空间 $\mathcal{H}$ 具有如下基底
\begin{equation*}
\left|n_{\pmb{k}_1}n_{\pmb{k}_2}\ldots\right\rangle
\end{equation*}
其中 $n_{\pmb{k}}=0,1,2,\ldots$ 是单粒子动量本征态 $|\pmb{k}\rangle$ 上的玻色子数。态 $|n_{\pmb{k}_1}n_{\pmb{k}_2}\ldots\rangle$ 的总能量为
\begin{equation*}
E=\sum_{\pmb{k}}\frac{\pmb{k}^2}{2m}\,n_{\pmb{k}}
\end{equation*}
总希尔伯特空间 $\mathcal{H}$ 也可以看作由矢量 $\pmb{k}\equiv(k_x,k_y,\ldots)=\frac{2\pi}{L}(n_x,n_y,\ldots)$ 所标记的一簇谐振子的希尔伯特空间，而 $n_{\pmb{k}}$ 正是标记谐振子 $\pmb{k}$ 能级的整数。从玻色子的总能量可以看出，这些谐振子是相互独立的，谐振子 $\pmb{k}$ 由下面的哈密顿量描述：
\begin{equation*}
H_{\pmb{k}}=\epsilon_{\pmb{k}}a_{\pmb{k}}^{\dagger}a_{\pmb{k}},\qquad \epsilon_{\pmb{k}}=\frac{\pmb{k}^2}{2m}\tag{3.1.1}
\end{equation*}
这样，在二次量子化中，自由玻色系统由一簇谐振子描述，其哈密顿量为
\begin{equation*}
H=\sum_{\pmb{k}}\epsilon_{\pmb{k}}a_{\pmb{k}}^{\dagger}a_{\pmb{k}}
\end{equation*}
总玻色子数算符为
\begin{equation*}
\hat{N}=\sum_{\pmb{k}}a_{\pmb{k}}^{\dagger}a_{\pmb{k}}
\end{equation*}
满足
\begin{equation*}
a_{\pmb{k}}|0\rangle=0
\end{equation*}
的谐振子基态是一个没有玻色子的态。第 $i$ 个玻色子携带动量 $\pmb{k}_i$ 的 $N$ 玻色子态由 $N$ 个 $a_{\pmb{k}}^{\dagger}$ 算符产生：
\begin{equation*}
\left|\pmb{k}_1\pmb{k}_2\ldots\pmb{k}_N\right\rangle=C(\pmb{k}_1,\ldots,\pmb{k}_N)\,a_{\pmb{k}_1}^{\dagger}\ldots a_{\pmb{k}_N}^{\dagger}|0\rangle
\end{equation*}
其中 $C(\pmb{k}_1,\ldots,\pmb{k}_N)$ 是归一化系数。只有当所有 $\pmb{k}_i$ 彼此互不相同时，才有 $C(\pmb{k}_1,\ldots,\pmb{k}_N)=1$。注意，$a_{\pmb{k}}^{\dagger}$（$a_{\pmb{k}}$）使玻色子数增加（减少）$1$。因此，$a_{\pmb{k}}^{\dagger}$ 称为玻色子的产生算符，而 $a_{\pmb{k}}$ 称为湮灭算符。产生算符与湮灭算符并不出现在玻色系统的第一次量子化描述中，因为在那种描述里我们只考虑玻色子数固定的态。

我们可以引入
\begin{equation*}
a(\pmb{x})=\sum_{\pmb{k}}a_{\pmb{k}}\mathcal{V}^{-1/2}\mathrm{e}^{\mathrm{i}\pmb{k}\pmb{x}}
\end{equation*}
或者，对无限系统，引入
\begin{equation*}
a(\pmb{x})=\int\frac{\mathrm{d}^d\pmb{k}}{(2\pi)^d}\,a_{\pmb{k}}\mathrm{e}^{\mathrm{i}\pmb{k}\pmb{x}}
\end{equation*}
可以看到，$a^{\dagger}(\pmb{x})$ 在 $\pmb{x}$ 处产生一个玻色子。态 $|\pmb{k}_1\pmb{k}_2\ldots\pmb{k}_N\rangle$ 的波函数现在可以如下算出：
\begin{equation*}
\psi(\pmb{x}_1,\ldots,\pmb{x}_N)=\langle\pmb{x}_1\ldots\pmb{x}_N|\pmb{k}_1\ldots\pmb{k}_N\rangle,\qquad |\pmb{x}_1\ldots\pmb{x}_N\rangle=a^{\dagger}(\pmb{x}_1)\ldots a^{\dagger}(\pmb{x}_N)|0\rangle
\end{equation*}
哈密顿量可以在实空间中写成
\begin{equation*}
H=\int \mathrm{d}^d\pmb{x}\;a^{\dagger}(\pmb{x})\left(-\frac{1}{2m}\frac{\mathrm{d}^2}{\mathrm{d}\pmb{x}^2}\right)a(\pmb{x})
\end{equation*}
算符 $a^{\dagger}(\pmb{x})a(\pmb{x})$ 统计位置 $\pmb{x}$ 处的玻色子数，因此玻色子密度算符为
\begin{equation*}
\rho(\pmb{x})=a^{\dagger}(\pmb{x})a(\pmb{x})
\end{equation*}
有了哈密顿量以及对应于各物理量的算符，我们就能通过计算物理算符的关联函数来计算玻色系统的物理性质。

谐振子描述使我们能够计算玻色系统的所有关联函数。特别地，对 $t>0$，
\begin{equation*}
\mathrm{i}G(\pmb{x}_f-\pmb{x}_i,t)=\left\langle a(\pmb{x}_f)\mathrm{e}^{-\mathrm{i}tH}a^{\dagger}(\pmb{x}_i)\right\rangle=\left\langle a(\pmb{x}_f,t)a^{\dagger}(\pmb{x}_i,0)\right\rangle
\end{equation*}
正是 2.1.1 节所讨论的自由粒子的传播子，而
\begin{equation*}
\mathrm{i}G(\pmb{k},t)=\left\langle a_{-\pmb{k}}\mathrm{e}^{-\mathrm{i}tH}a_{\pmb{k}}^{\dagger}\right\rangle=\mathrm{e}^{-\mathrm{i}\epsilon_{\pmb{k}}t}
\end{equation*}
是相应的动量空间中的传播子。

要计算多个玻色算符的关联，例如 $\rho$--$\rho$ 关联，我们需要使用 Wick 定理。Wick 定理对于做正规排序（normal ordering）和计算关联都非常有用。

\textbf{定理.}（Wick 定理）设 $O_i$ 是 $a_{\pmb{k}}$ 与 $a_{\pmb{k}}^{\dagger}$ 的线性组合：$O_i=\int \mathrm{d}k\,\left(u_i(k)a_{\pmb{k}}+v_i(k)a_{\pmb{k}}^{\dagger}\right)$。（注意，对自由玻色系统，$a(t)=U^{\dagger}(t,0)aU(t,0)$ 就是这样一个算符。）设 $W=\prod_{i=1}^{N}O_i$，$W_{i_1,i_2}=\prod_{i=1,i\neq i_1,i\neq i_2}^{N}O_i$，等等。设 $:W:$ 是算符 $W$ 的正规排序形式（即把所有 $a_{\pmb{k}}^{\dagger}$ 都放到 $a_{\pmb{k}}$ 的左边），并设 $|0\rangle$ 是被所有 $a_{\pmb{k}}$ 湮灭的态（即对一切 $\pmb{k}$ 有 $a_{\pmb{k}}|0\rangle=0$）。那么 Wick 定理说：
\begin{align*}
W={}&:W:+\sum_{(i_1,j_1)}:W_{i_1,j_1}:\,\langle 0|O_{i_1}O_{j_1}|0\rangle\\
&+\sum_{\substack{(i_1,j_1),(i_2,j_2)\\ (i_1,j_1)\neq(i_2,j_2)}}:W_{i_1,j_1,i_2,j_2}:\,\langle 0|O_{i_1}O_{j_1}|0\rangle\langle 0|O_{i_2}O_{j_2}|0\rangle+\cdots\tag{3.1.2}
\end{align*}
其中 $(i_1,j_1)$ 是满足 $i_1<j_1$ 的有序对。

为了展示 Wick 定理的威力，我们考虑如下关联：
\begin{equation*}
\langle 0|A_1A_2A_3A_4|0\rangle
\end{equation*}
其中诸 $A_i$ 是 $a$ 与 $a^{\dagger}$ 的线性组合。由于 $\langle 0|:A_iA_j:|0\rangle=0$，Wick 定理给出
\begin{align*}
&\langle 0|A_1A_2A_3A_4|0\rangle\\
&\quad=\langle 0|A_1A_2|0\rangle\langle 0|A_3A_4|0\rangle+\langle 0|A_1A_3|0\rangle\langle 0|A_2A_4|0\rangle+\langle 0|A_1A_4|0\rangle\langle 0|A_2A_3|0\rangle
\end{align*}
这样，四算符关联就可以用两算符关联来表达。

\subsection*{问题 3.1.1}

Wick 定理：
\begin{enumerate}
\item 对算符 $O=aa\,a^{\dagger}a^{\dagger}$ 进行正规排序，并证明 Wick 定理给出正确的结果。
\item 用 Wick 定理计算 $\langle 0|aa^{\dagger}aa^{\dagger}|0\rangle$。
\end{enumerate}

\subsection*{问题 3.1.2}

考虑一个含 $N$ 个玻色子的一维自由玻色系统，令 $|\Phi_0\rangle$ 为其基态。
\begin{enumerate}
\item 对 $N=0$ 与有限的 $N$，计算编时传播子
\begin{equation*}
\mathrm{i}G(x,t)=\langle\Phi_0|T(a(x,t)a^{\dagger}(0,0))|\Phi_0\rangle
\end{equation*}
其中 $a(x,t)=\mathrm{e}^{\mathrm{i}Ht}a(x)\mathrm{e}^{-\mathrm{i}Ht}$。证明 $\mathrm{i}G(x,0^{+})-\mathrm{i}G(x,0^{-})=[a(x),a^{\dagger}(0)]=\delta(x)$。证明：对有限的 $N$，在 $x\to\infty$ 极限下有 $\mathrm{i}G(x,t)\neq0$，这表明存在（非对角）长程序（off-diagonal long-range order）。非对角长程序只存在于玻色凝聚态之中。
\item 计算编时密度关联函数
\begin{equation*}
\langle\Phi_0|T(\rho(x,t)\rho(0,0))|\Phi_0\rangle
\end{equation*}
证明 $a(x,t)$ 是 $a_{\pmb{k}}$ 与 $a_{\pmb{k}}^{\dagger}$ 的线性组合，因而可以使用 Wick 定理。（提示：你可能需要先把 $a_0$ 算符分离出来。）
\end{enumerate}

\section{超流的平均场理论}

\begin{keypoints}
\item 平均场理论这样得到：忽略一些算符的量子涨落，并用 c 数（c-number，通常取为这些算符的平均值）代替它们。
\item 平均场基态可以看作一个尝试波函数（trial wave function）。平均场理论所给出的激发谱可能在定性上就是错的。在使用平均场激发谱之前，需要先用其他手段论证其正确性。
\end{keypoints}

零温下，自由的 $N$ 玻色系统处于其基态
\begin{equation*}
\left|\Phi_0\right\rangle=(N!)^{-1/2}\left(a_0^{\dagger}\right)^N|0\rangle
\end{equation*}
所有玻色子都处于 $\pmb{k}=0$ 的态。这个态称为玻色凝聚态（boson condensed state）。自由玻色系统的玻色凝聚态具有许多特殊性质，这些性质只对无相互作用的玻色子才成立。要理解真实玻色系统的零温态，我们需要研究由如下哈密顿量描述的相互作用玻色系统：
\begin{equation*}
H=\sum_{\pmb{k}}(\epsilon_{\pmb{k}}-\mu)a_{\pmb{k}}^{\dagger}a_{\pmb{k}}+\int \mathrm{d}^d\pmb{x}\,\mathrm{d}^d\pmb{x}'\,\frac{1}{2}\rho(\pmb{x})V(\pmb{x}-\pmb{x}')\rho(\pmb{x}')
\end{equation*}
其中 $V(\pmb{x})$ 是密度--密度相互作用，$\mu$ 是化学势（chemical potential）。$\mu$ 的取值使得系统中平均有 $N$ 个玻色子。由于包含了化学势项 $-\mu N$，$H$ 是热势（thermal potential）的算符，而不是能量的算符。在动量空间中，我们有
\begin{align*}
H&=\sum_{\pmb{k}}(\epsilon_{\pmb{k}}-\mu)a_{\pmb{k}}^{\dagger}a_{\pmb{k}}+\mathcal{V}\sum_{\pmb{q}}\frac{1}{2}\rho_{-\pmb{q}}V_{\pmb{q}}\rho_{\pmb{q}}\\
&=\sum_{\pmb{k}}(\epsilon_{\pmb{k}}-\mu)a_{\pmb{k}}^{\dagger}a_{\pmb{k}}+\mathcal{V}^{-1}\frac{1}{2}\sum_{\pmb{k},\pmb{k}',\pmb{q}}\left(a_{\pmb{k}'-\pmb{q}}^{\dagger}a_{\pmb{k}'}\right)V_{\pmb{q}}\left(a_{\pmb{k}+\pmb{q}}^{\dagger}a_{\pmb{k}}\right)\\
&=\sum_{\pmb{k}}(\epsilon_{\pmb{k}}-\mu+V(0))a_{\pmb{k}}^{\dagger}a_{\pmb{k}}+\mathcal{V}^{-1}\frac{1}{2}\sum_{\pmb{k},\pmb{k}',\pmb{q}}V_{\pmb{q}}\,a_{\pmb{k}'-\pmb{q}}^{\dagger}a_{\pmb{k}+\pmb{q}}^{\dagger}a_{\pmb{k}'}a_{\pmb{k}}\tag{3.2.1}
\end{align*}
其中
\begin{equation*}
a_{\pmb{k}}=\int \mathrm{d}\pmb{x}\,a(\pmb{x})\mathcal{V}^{-1/2}\mathrm{e}^{\mathrm{i}\pmb{k}\cdot\pmb{x}},\qquad \rho_{\pmb{k}}=\int \mathrm{d}\pmb{x}\,\rho(\pmb{x})\mathrm{e}^{\mathrm{i}\pmb{k}\cdot\pmb{x}},\qquad V_{\pmb{k}}=\int \mathrm{d}\pmb{x}\,V(\pmb{x})\mathrm{e}^{\mathrm{i}\pmb{k}\cdot\pmb{x}}
\end{equation*}
在式 (3.2.1) 的最后一行里，相互作用项已被写成正规排序的形式，即 $a^{\dagger}$ 总是出现在 $a$ 的左边。此外，由于 $V(\pmb{x})$ 是对称的，即 $V(\pmb{x})=V(-\pmb{x})$，$V_{\pmb{k}}$ 为实数，且 $V_{\pmb{k}}=V_{-\pmb{k}}$。下面我们将把 $V(0)$ 吸收进 $\mu$ 并略去 $V(0)$。这个相互作用哈密顿量在 $a_{\pmb{k}}$ 上不是二次型的，很难求解。

下面我们将用平均场方法来求一个近似基态。平均场理论的基本想法是：找出一些量子涨落较弱的算符，用相应的经典值（一个 c 数）去代替它们。但愿这一替换能把完整的哈密顿量约化成二次型的。

基于对自由玻色子基态的知识，我们可以假设：在相互作用玻色系统的基态中，有宏观数目的粒子占据 $\pmb{k}=0$ 态：
\begin{equation*}
\langle\Phi_0|a_0^{\dagger}a_0|\Phi_0\rangle=N_0
\end{equation*}
而其他态的占据数，即 $n_{\pmb{k}}=\langle\Phi_0|a_{\pmb{k}}^{\dagger}a_{\pmb{k}}|\Phi_0\rangle$，都很小。我们还注意到
\begin{align*}
&\langle\Phi_0|a_0^{\dagger}a_0a_0^{\dagger}a_0|\Phi_0\rangle=N_0^2\\
&\langle\Phi_0|a_0^{\dagger}a_0^{\dagger}a_0a_0|\Phi_0\rangle=N_0(N_0-1)\approx N_0^2
\end{align*}
这样，精确到 $N_0^2$ 阶，$a_0$ 与 $a_0^{\dagger}$ 看上去彼此对易，表现得像普通的数（c 数）。这就是我们说算符 $a_0$ 的量子涨落很弱时的含义。在这种情形下，我们可以用它的经典值代替 $a_0$：
\begin{equation*}
a_0=a_0^{\dagger}=\sqrt{N_0}
\end{equation*}
相互作用哈密顿量这时可以写成
\begin{align*}
H_{mean}&=\sum_{\pmb{k}}(\epsilon_{\pmb{k}}-\mu)a_{\pmb{k}}^{\dagger}a_{\pmb{k}}+O\left((a_{\pmb{k}\neq0})^3\right)+\frac{\mathcal{V}}{2}\rho_0^2V_0\\
&\quad+\frac{\rho_0}{2}\sum_{\pmb{k}\neq0}\left(2a_{\pmb{k}}^{\dagger}a_{\pmb{k}}V_0+2a_{\pmb{k}}^{\dagger}a_{\pmb{k}}V_{\pmb{k}}+V_{\pmb{k}}a_{-\pmb{k}}a_{\pmb{k}}+V_{\pmb{k}}a_{-\pmb{k}}^{\dagger}a_{\pmb{k}}^{\dagger}\right)\\
&=\sum_{\pmb{k}}(\epsilon_{\pmb{k}}-\mu'_{\pmb{k}})a_{\pmb{k}}^{\dagger}a_{\pmb{k}}-\frac{1}{2}\rho_0^2V_0+\rho_0\frac{1}{2}\sum_{\pmb{k}\neq0}V_{\pmb{k}}\left(a_{-\pmb{k}}a_{\pmb{k}}+a_{-\pmb{k}}^{\dagger}a_{\pmb{k}}^{\dagger}\right)\tag{3.2.2}
\end{align*}
其中 $\mu'_{\pmb{k}}=\mu-\rho_0V_0-\frac{1}{2}\rho_0(V_{\pmb{k}}+V_{-\pmb{k}})$，$\rho_0=N_0/\mathcal{V}$，并且我们忽略了 $(a_{\pmb{k}\neq0})^3$ 以及更高阶的项。经过两次近似——第一步用 c 数代替 $a_0$（这在 $N\to\infty$ 极限下是准确的），第二步略去 $a_{\pmb{k}}^3$ 项（这在 $V\to0$ 或 $\rho\to0$ 极限下是准确的）——我们便得到一个二次型的平均场哈密顿量。

由于出现了 $a_{-\pmb{k}}a_{\pmb{k}}$ 项，显然，所有 $N$ 个玻色子都处于 $\pmb{k}=0$ 态的态不再是相互作用玻色系统的基态。为了找到新的（平均场）基态，我们对 $\pmb{k}\neq0$ 引入
\begin{equation*}
\alpha_{\pmb{k}}=u_{\pmb{k}}a_{\pmb{k}}+v_{\pmb{k}}a_{-\pmb{k}}^{\dagger}
\end{equation*}
我们这样选取 $u_{\pmb{k}}$ 与 $v_{\pmb{k}}$，使得 $H_{mean}$ 具有 $\sum_{\pmb{k}}\alpha_{\pmb{k}}^{\dagger}\alpha_{\pmb{k}}E_{\pmb{k}}+\text{常数}$ 的形式，且 $\alpha_{\pmb{k}}$ 与 $\alpha_{\pmb{k}}^{\dagger}$ 满足同样的玻色子对易关系
\begin{equation*}
\left[\alpha_{\pmb{k}},\alpha_{\pmb{k}'}^{\dagger}\right]=\delta_{\pmb{k},\pmb{k}'}
\end{equation*}
后一条件要求 $|u_{\pmb{k}}^2|-|v_{\pmb{k}}|^2=1$。可以证明，如下选取满足我们的要求：
\begin{align*}
E_{\pmb{k}}&=\sqrt{(\epsilon_{\pmb{k}}-\mu'_{\pmb{k}})^2-\rho_0^2V_{\pmb{k}}^2}\tag{3.2.3}\\
u_{\pmb{k}}&=\sqrt{\frac{\epsilon_{\pmb{k}}-\mu'_{\pmb{k}}}{2E_{\pmb{k}}}+\frac{1}{2}}\tag{3.2.4}\\
v_{\pmb{k}}&=\frac{V_{\pmb{k}}}{|V_{\pmb{k}}|}\sqrt{\frac{\epsilon_{\pmb{k}}-\mu'_{\pmb{k}}}{2E_{\pmb{k}}}-\frac{1}{2}}\tag{3.2.5}
\end{align*}
并且
\begin{equation*}
\sum_{\pmb{k}\neq0}\left((\epsilon_{\pmb{k}}-\mu'_{\pmb{k}})a_{\pmb{k}}^{\dagger}a_{\pmb{k}}+\frac{V_{\pmb{k}}\rho_0}{2}\left(a_{-\pmb{k}}a_{\pmb{k}}+\text{h.c.}\right)\right)=\sum_{\pmb{k}\neq0}E_{\pmb{k}}\alpha_{\pmb{k}}^{\dagger}\alpha_{\pmb{k}}-\sum_{\pmb{k}\neq0}E_{\pmb{k}}|v_{\pmb{k}}|^2
\end{equation*}
把 $\pmb{k}=0$ 的部分也包括进来，我们得到
\begin{align*}
H_{mean}&=\sum_{\pmb{k}\neq0}E_{\pmb{k}}\alpha_{\pmb{k}}^{\dagger}\alpha_{\pmb{k}}+\Omega_g\\
\Omega_g&=\mathcal{V}\left(-\mu\rho_0+\frac{1}{2}\rho_0^2V_0\right)-\sum_{\pmb{k}\neq0}\frac{1}{2}(\epsilon_{\pmb{k}}-\mu'_{\pmb{k}}-E_{\pmb{k}})
\end{align*}
平均场基态由
\begin{equation*}
\alpha_{\pmb{k}}|\Phi_{mean}\rangle=0
\end{equation*}
给出，而 $\Omega_g$ 就是平均场基态能量，或者更准确地说，是平均场基态的热势。在平均场理论中，基态之上的激发由一簇谐振子 $\alpha_{\pmb{k}}$ 描述，它们对应于玻色流体中的波。波的色散关系（dispersion relation）由 $E_{\pmb{k}}$ 给出。

还有两个参量 $N_0$ 与 $\mu$ 待定。激发谱
\begin{equation*}
E_{\pmb{k}}=\sqrt{(\epsilon_{\pmb{k}}-\mu+\rho_0V_0+\rho_0V_{\pmb{k}})^2-(\rho_0V_{\pmb{k}})^2}
\end{equation*}
以一种敏感的方式依赖于 $\mu$。若 $\mu=\rho_0V_0$，则 $E_0=0$，激发是无能隙的（这里我们假设了 $\epsilon_0=0$）。若 $\mu<\rho_0V_0$，则激发具有有限的能隙。若 $\mu>\rho_0V_0$，则 $E_0$ 是虚数，这暗示着一种不稳定性。

首先，$N_0$ 的值通过求热势 $\Omega_g$ 的极小（保持 $\mu$ 固定）得到：
\begin{equation*}
\frac{\partial\Omega_g}{\partial N_0}=0\tag{3.2.6}
\end{equation*}
求得 $N_0=N_0(\mu)$ 之后，注意到
\begin{equation*}
a_{\pmb{k}}=u_{\pmb{k}}^{*}\alpha_{\pmb{k}}-v_{\pmb{k}}^{*}\alpha_{-\pmb{k}}^{\dagger}
\end{equation*}
我们发现
\begin{equation*}
N=\langle\Phi_{mean}|\sum_{\pmb{k}}a_{\pmb{k}}^{\dagger}a_{\pmb{k}}|\Phi_{mean}\rangle=N_0(\mu)+\sum_{\pmb{k}}|v_{\pmb{k}}|^2\tag{3.2.7}
\end{equation*}
上式确定了给出总共 $N$ 个玻色子的 $\mu=\mu(N)$ 的值。

让我们尝试在低密度或弱耦合极限下进行上述计算。保留 $V_{\pmb{k}}$ 的线性阶项（注意 $\epsilon_{\pmb{k}}-\mu'_{\pmb{k}}-E_{\pmb{k}}=O(V_{\pmb{k}}^2)$），我们得到
\begin{equation*}
\Omega_g=\mathcal{V}\left(-\mu\rho_0+\frac{1}{2}\rho_0^2V_0\right)
\end{equation*}
我们看到，为使 $\Omega_g$ 极小，要求 $\mu<0$ 时 $N_0=0$，而 $\mu>0$ 时 $N_0=\mathcal{V}\mu/V_0$。一旦知道了 $N_0$，玻色子总数就可以通过式 (3.2.7) 算出。如果已知 $N$，则 $\mu$ 由 $N=\mathcal{V}\mu/V_0+\sum_{\pmb{k}}|v_{\pmb{k}}|^2$ 确定。

当玻色子密度非零时，我们求得 $\mu=\rho_0V_0>0$，激发谱
\begin{equation*}
E_{\pmb{k}}=\sqrt{\epsilon_{\pmb{k}}\left(\epsilon_{\pmb{k}}+2\rho_0V_{\pmb{k}}\right)}\tag{3.2.8}
\end{equation*}
对小 $\pmb{k}$ 是无能隙且线性的，不论 $N$ 与 $V_{\pmb{k}}$ 的取值如何：
\begin{equation*}
E_{\pmb{k}}=v|\pmb{k}|,\qquad v^2=\frac{\rho_0V_0}{m}-\frac{\rho_0^2}{2}\tag{3.2.9}
\end{equation*}
（译注：按式 (3.2.8) 在小 $\pmb{k}$ 下展开，应得 $v^2=\rho_0V_0/m$；原书式中的 $-\rho_0^2/2$ 一项疑为排印之误。）

这个结果非常合理，因为 $\alpha_{\pmb{k}}^{\dagger}$ 产生的激发应当对应于可压缩玻色流体中的密度波（即 Nambu--Goldstone 模），因而理应总具有无能隙的线性色散。这种线性无能隙激发也称为声子（phonon）。

如果我们就此止步，那么一切都会很好。可以说，我们得到了一个关于相互作用玻色系统基态与激发的平均场理论，它给出合理的结果。然而，如果我们想做得更好，事情反而可能变糟。在平均场理论中，能隙的消失是由于一个"奇迹"——$\mu$ 恰好等于 $\rho_0V_0$，从而使 $E_{\pmb{k}}$ 中的常数项相消。如果我们保留 $\Omega_g$ 中出现的 $V_{\pmb{k}}$ 的更高阶项，那么更精确的 $\mu$ 可能就不再等于 $\rho_0V_0$。在这种情形下，我们可能得到能隙不为零的不合理结果。当然，当我们在 $\Omega_g$ 中计入更高阶项时，我们应当同时在 $E_{\pmb{k}}$ 中计入更高阶项，这时仍可能发生相消而给出零能隙。

%FIG{3.1}: {声子激发位于 $\pmb{k}=0$ 附近，旋子激发位于有限 $\pmb{k}$ 处的局域极小附近。图中还标出了超流流动的临界速度（参见 3.7.3 节）。}

然而，有一点是清楚的。平均场理论并不保证零能隙。要得到零能隙，我们只能依靠仔细的计算并祈盼奇迹。但是，从物理的角度看，无能隙激发的存在是一条普适原理，它不依赖于玻色系统的任何细节。无论相互作用取何种形式，自由空间中的玻色系统总是可压缩的，并且总含有无能隙的密度波。我们看到，平均场理论并没有把握住无能隙激发背后的原理。平均场理论的这个问题并不只出现在玻色系统中。一般而言，所有平均场理论都有这样的毛病：原理论的某些普遍原理没有被平均场理论捕捉到。在用平均场理论研究激发谱时，必须牢记这一点，并格外小心。

在真实的超流体中，激发谱具有图 3.1 所示的形状。位于有限 $\pmb{k}$ 处极小附近的激发称为旋子（roton）。我们注意到，在适当的 $\pmb{k}$ 处，旋子的群速度可以为零。

\subsection*{问题 3.2.1}

对于一个密度为 $\rho=N/\mathcal{V}$ 的一维相互作用玻色系统，确定 $N_0$ 与 $\mu$。证明一维的相互作用玻色子即使在零温下也不能凝聚。（注意，无相互作用的玻色子在零温下可以在一维凝聚。）

\subsection*{问题 3.2.2}

\begin{enumerate}
\item 找出一个 $V_{\pmb{k}}$，使得 $E_{\pmb{k}}$ 中的旋子极小降到零以下。在实空间中画出这个 $V(\pmb{x})$。
\item 当旋子极小果真降到零以下时，超流态会发生什么？
\end{enumerate}

\subsection*{问题 3.2.3}

自旋链的平均场理论。考虑自旋链
\begin{equation*}
H=J\sum_i S_iS_{i+1}
\end{equation*}
其中 $J<0$，$S$ 是携带自旋 $S$ 的自旋算符。当 $S\gg0$ 时，$S_i$ 的量子涨落很弱，我们可以把一个 $S$ 用其平均值 $s=\langle S_i\rangle$（假设它不依赖 $i$）代替，从而得到平均场哈密顿量。
\begin{enumerate}
\item 写出平均场哈密顿量，并求出平均场基态。
\item 用平均场哈密顿量计算 $\langle S_i\rangle$，从而确定 $s$ 的值。
\item 用平均场近似求出基态能量。
\item 求出平均场基态之上各激发的能量。
\end{enumerate}

我们由此看到，平均场理论完全无法重现无能隙的自旋波激发。2.4 节所讨论的路径积分方法及其经典近似则会给出好得多的结果。

\section{相互作用玻色系统的路径积分方法}

\subsection{相互作用玻色系统的路径积分表示}

如 3.1 节所述，自由玻色系统可以用一簇独立的谐振子来描述，其哈密顿量为
\begin{equation*}
H_0=\sum_{\pmb{k}}(\epsilon_{\pmb{k}}-\mu)a_{\pmb{k}}^{\dagger}a_{\pmb{k}}
\end{equation*}
这些谐振子（以及自由玻色系统）容许一个（相干态）路径积分表示
\begin{equation*}
\int\mathcal{D}^2\left[a_{\pmb{k}}(t)\right]\,\mathrm{e}^{\mathrm{i}\int \mathrm{d}t\,\sum_{\pmb{k}}\left[\mathrm{i}\frac{1}{2}(a_{\pmb{k}}^{*}\dot{a}_{\pmb{k}}-a_{\pmb{k}}\dot{a}_{\pmb{k}}^{*})-(\epsilon_{\pmb{k}}-\mu)a_{\pmb{k}}^{*}a_{\pmb{k}}\right]}
\end{equation*}
在实空间中，我们有
\begin{equation*}
\int\mathcal{D}^2a\;\mathrm{e}^{\mathrm{i}\int \mathrm{d}^d\pmb{x}\,\mathrm{d}t\,\left[\mathrm{i}\frac{1}{2}\left(a^{*}(\pmb{x},t)\partial_t a(\pmb{x},t)-a(\pmb{x},t)\partial_t a^{*}(\pmb{x},t)\right)-\frac{1}{2m}\partial_{\pmb{x}}a^{*}(\pmb{x},t)\partial_{\pmb{x}}a(\pmb{x},t)-\mu a^{*}(\pmb{x},t)a(\pmb{x},t)\right]}
\end{equation*}
现在我们可以很容易地把相互作用包括进来，只需在哈密顿量与路径积分中加上一项
\begin{equation*}
\int \mathrm{d}^d\pmb{x}\,\mathrm{d}^d\pmb{y}\;\frac{1}{2}a^{*}(\pmb{x},t)a(\pmb{x},t)V(\pmb{x}-\pmb{y})a^{*}(\pmb{y},t)a(\pmb{y},t)
\end{equation*}
这里我们将处理一个更简单的相互作用 $V(\pmb{x})=V_0\delta(\pmb{x})$，其傅里叶分量为 $V_{\pmb{k}}=V_0$。相互作用玻色子的路径积分于是具有形式
\begin{equation*}
\int\mathcal{D}^2\left[\varphi(\pmb{x},t)\right]\,\mathrm{e}^{\mathrm{i}\int \mathrm{d}^d\pmb{x}\,\mathrm{d}t\,\left[\mathrm{i}\frac{1}{2}(\varphi^{*}\partial_t\varphi-\varphi\partial_t\varphi^{*})-\frac{1}{2m}\partial_{\pmb{x}}\varphi^{*}\partial_{\pmb{x}}\varphi+\mu|\varphi|^2-\frac{V_0}{2}|\varphi|^4\right]}
\end{equation*}
其中，按照惯例，我们已把 $a(\pmb{x},t)$ 改记为 $\varphi(\pmb{x},t)$。

%FIG{3.2}: {热势 $\Omega(\mu)$ 的二阶导数出现不连续，这标志着一次二级相变。}

有了路径积分表示，我们就可以把这个量子相互作用玻色系统当作一个经典场论来研究其物理性质，该经典场论由作用量
\begin{equation*}
S=\int \mathrm{d}^d\pmb{x}\,\mathrm{d}t\,\left[\mathrm{i}\frac{1}{2}(\varphi^{*}\partial_t\varphi-\varphi\partial_t\varphi^{*})-\frac{1}{2m}\partial_{\pmb{x}}\varphi^{*}\partial_{\pmb{x}}\varphi+\mu|\varphi|^2-\frac{V_0}{2}|\varphi|^4\right]\tag{3.3.1}
\end{equation*}
描述。这里我们做的是半经典近似（而不是 3.2 节所做的平均场近似），其领头项就是经典理论。这个经典系统的能量（或者，更准确地说，热势）为
\begin{equation*}
\Omega=\int \mathrm{d}^d\pmb{x}\,\left[\frac{1}{2m}\partial_{\pmb{x}}\varphi^{*}\partial_{\pmb{x}}\varphi+\mu|\varphi|^2+\frac{V_0}{2}|\varphi|^4\right]\tag{3.3.2}
\end{equation*}

\subsection{相变与自发对称性破缺}

\begin{keypoints}
\item 零温相变发生在基态能量的非解析点处。
\item 当一个非对称的基态从对称系统中涌现时，就发生了自发对称性破缺。
\item 连续相变、自发对称性破缺与长程序密切相关。这正是 Landau 关于相与相变的对称性破缺理论的核心。
\end{keypoints}

经典基态是平移不变的，由一个复常数 $\varphi_0$ 刻画：$\varphi(\pmb{x},t)=\varphi_0$。经典基态的能量密度（或者，更准确地讲，是热势的密度）为

% ------------------------------------------------------------
% 块边界备注：
%   终点：书页 78（p093）末句 "The energy density (or, more precisely, the
%   density of the thermal potential) of the classical ground state is"
%   跨页延续到书页 79（p094）顶部。p094 顶部残句 "density of the thermal
%   potential) ..." 的译文（"密度）为"）已按任务约定写入本文件顶部
%   %--A-TAIL-- 区。ch3_b 请勿重复翻译该残句，正文请从 p094 的无编号式
%   Ω_0/𝒱 = −μ|φ_0|² + (V_0/2)|φ_0|⁴ 起译（其后为 "By minimizing the
%   energy..." 一段）。
%   另：本块含 §3.3 与 §3.3.1 的开头、§3.3.2 的标题＋要点＋首句（未完）。
% ------------------------------------------------------------
% 新增术语（ch3_a，待并入术语表"新增术语"区）：
%   first quantization 第一次量子化
%   normal ordering 正规排序
%   c-number c 数
%   boson condensed state 玻色凝聚态
%   chemical potential 化学势
%   dispersion relation 色散关系
%   off-diagonal long-range order 非对角长程序
%   trial wave function 尝试波函数
%   phonon 声子
%   roton 旋子
%   critical velocity 临界速度
%   group velocity 群速度
%   density wave 密度波
%   Nambu–Goldstone modes Nambu–Goldstone 模
%   second-order phase transition 二级相变
%   classical field theory 经典场论
%   Fourier component 傅里叶分量
%   spin chain 自旋链
%   spin wave 自旋波
%   weak-coupling limit 弱耦合极限
%   low-density limit 低密度极限
